\section{বিশ্বগ্রাম প্রতিষ্ঠার উপাদান ও তথ্য ও যোগাযোগ প্রযুক্তির অবদান}

\begin{learningbox}
এই টপিক শেষে শিক্ষার্থী পারবে:
\begin{itemize}
\item বিশ্বগ্রাম প্রতিষ্ঠার প্রধান উপাদানগুলো চিহ্নিত করতে।
\item হার্ডওয়্যার, সফটওয়্যার, ডেটা, নেটওয়ার্ক ও দক্ষতার ভূমিকা ব্যাখ্যা করতে।
\item শিক্ষা, চিকিৎসা, ব্যবসা, গবেষণা ও কর্মসংস্থানে ICT-এর অবদান লিখতে।
\end{itemize}
\end{learningbox}

\subsection{বিশ্বগ্রাম প্রতিষ্ঠার প্রধান উপাদান}
\subsection{হার্ডওয়্যার}
\subsection{সফটওয়্যার}
\subsection{ডেটা ও তথ্যভান্ডার}
\subsection{নেটওয়ার্ক ও ইন্টারনেট}
\subsection{মানুষের জ্ঞান ও দক্ষতা}
\subsection{শিক্ষায় তথ্য ও যোগাযোগ প্রযুক্তির অবদান}
\subsection{চিকিৎসায় তথ্য ও যোগাযোগ প্রযুক্তির অবদান}
\subsection{ব্যবসা-বাণিজ্যে তথ্য ও যোগাযোগ প্রযুক্তির অবদান}
\subsection{গবেষণায় তথ্য ও যোগাযোগ প্রযুক্তির অবদান}
\subsection{কর্মসংস্থানে তথ্য ও যোগাযোগ প্রযুক্তির অবদান}

\begin{recapbox}
\begin{itemize}
\item বিশ্বগ্রাম শুধু internet নয়; hardware, software, data, network ও skilled user মিলেই এটি তৈরি হয়।
\item ICT মানুষের যোগাযোগ, সেবা গ্রহণ, শেখা, কাজ ও ব্যবসার পদ্ধতিকে দ্রুত করেছে।
\item এই topic থেকে বোর্ড পরীক্ষায় উপাদান, ICT-এর অবদান ও উদাহরণভিত্তিক প্রশ্ন আসতে পারে।
\end{itemize}
\end{recapbox}
